\chapter{Эксперимент: как это проверять}
\label{ch:experiment}

\section{Эрзац-граф лучше живой Рин}

Тестировать архитектуру памяти на готовом агенте — долго и слепо: сначала собери всё,
а потом окажется, что с памятью куча неоптимальностей, найденных \emph{после} финальной
сборки. Нужен управляемый стенд — эрзац-граф с retrieval.

\begin{principle}[Ground truth важнее реализма]
Суррогатный граф — не бедный заменитель, а для отладки механик \emph{лучше} живой Рин,
потому что синтетика даёт \emph{ground truth}. С настоящей Рин ты не знаешь, что она
«должна» помнить, значит не измеришь, правильно ли отработали fusion, скраб, пропагация.
С подсаженной синтетикой истину задаёшь ты — и меряешь точность/полноту починки,
сходимость уверенности, нашёл ли скраб подсаженную ложь.
\end{principle}

Это тот же принцип, что контролируемая оценка в DTI: управляемый синтетический тест
изолирует эффект \emph{до} грязных реальных данных. Слабые места находятся \emph{до}
финальной сборки.

\section{Три яруса}

\paragraph{Ярус 1 — синтетическая топология (начинать отсюда).} Графы с известной
структурой, чтобы юнит-тестить и диагностику, и механики против теории: стохастическая
блочная модель (SBM — сообщества с регулируемым мостом, для over-squashing/проводимости),
Барабаши–Альберт (scale-free хабы — для «гравитационных ям» и концентрации $\pi$),
Уоттс–Строгац (small-world), деревья против экспандеров (крайности кривизны). Семантики
ноль, зато полная контролируемость: подсаживаешь ложное-правдоподобное убеждение и
противоречащее свидетельство в угол с малым $\pi$ и меряешь, находит ли телепорт-скраб.

\paragraph{Ярус 2 — семантический граф из готовых баз.} Чтобы у fusion/retrieval был
смысл. Разбор баз — раздел~\ref{sec:datasets}.

\paragraph{Ярус 3 — LLM-генерённый «поток жизни».} Полная петля с реалистичным
содержанием и контролируемой истиной: языковая модель генерит синтетическую
автобиографию/эпизоды/диалоги фейк-персоны с подсаженными фактами, противоречиями,
эволюцией убеждений. Прямой прообраз — \emph{Generative Agents} (Park и др.): поток
памяти $+$ retrieval (свежесть$\times$важность$\times$релевантность) $+$ рефлексия (это
и есть fusion). «Ползучую веру» подсаживаешь вручную и меряешь, лечится ли.

\section{Разбор баз знаний}
\label{sec:datasets}

\begin{center}
\begin{tabular}{p{0.16\textwidth} p{0.30\textwidth} p{0.42\textwidth}}
\toprule
\textbf{База} & \textbf{Что даёт} & \textbf{Под какую механику} \\
\midrule
ConceptNet & ассоциативно-таксономические связи, взвешенные & spreading activation,
ассоциативный fusion; веса $\approx$ мягкая уверенность \\
ATOMIC/10x & причинно-социальный смысл событий (намерение/реакция/эффект) &
эпизодика, сцепка с настроением, fusion-как-вывод \\
Wikidata & факты $+$ схема (время, references, ранг) & \emph{схему} красть под
темпоральную валидность, провенанс, ревизию; содержание — нет \\
\bottomrule
\end{tabular}
\end{center}

Не «какой из них», а «какой под какую механику». На старте — ConceptNet (чистый граф
под первую механику); Wikidata-\emph{схему} — когда дойдёшь до темпоральных слоёв и
ревизии. Порознь одного мало: ConceptNet без ATOMIC — мир без психологии, ATOMIC без
ConceptNet — реакции без предметного мира.

\section{Эмбриональный граф и работа стыка}

Объединить ConceptNet и ATOMIC в \emph{эмбриональный} граф и доращивать фьюжнами —
канонично: базовая когнитивная матрица, в которую новые факты встраиваются рёбрами.

\begin{pitfall}[«Форматы отличаются только на диске» — недооценка]
Внутри графовой базы это одна структура, и математика обхода от типа ребра не меняется.
Но настоящая работа не в формате, а в \emph{интерфейсе}: узлы ConceptNet — концепты,
узлы ATOMIC — события с плейсхолдерами \textsc{PersonX/Y}. Связать «разбить кружку» с
концептом «кружка» и заземлить \textsc{PersonX/Y} на реальных собеседников — это
entity/event linking (через Numberbatch, глава~\ref{ch:voice}), нетривиальный шаг, а
не объединение таблиц.
\end{pitfall}

\section{Батарея экспериментов}

Каждый тест — одна механика плюс метрика; всё гоняется уже на Ярусе 1.

\begin{center}
\begin{tabular}{p{0.05\textwidth} p{0.44\textwidth} p{0.42\textwidth}}
\toprule
\# & \textbf{Гипотеза} & \textbf{Установка и метрика} \\
\midrule
1 & $\pi$ концентрируется на хабах & BA-граф; распределение посещений против степени \\
2 & узкое место душит cross-cluster fusion & два SBM-сообщества; частота fusion против
ширины моста $m$ \\
3 & телепорт-скраб находит ложь в углу с малым $\pi$ & подсадка; precision/recall и
время-до-нахождения против $\alpha$ \\
4 & уверенность пересглаживается за много циклов & дисперсия $\conf$ во времени; тест
якоря-остатка \\
5 & руминация и её гашение & mood-congruent $+$ инерция; автокорреляция $\mood$ с
setpoint и без \\
\bottomrule
\end{tabular}
\end{center}

Скелет стенда (build $\to$ plant $\to$ operator $\to$ metric) — в
приложении~\ref{ch:algorithms}.

\section{Микроэксперименты: наблюдать физику потока}

Помимо «прошёл/не прошёл»-тестов есть класс более тонких наблюдений, вырастающих
прямо из главы~\ref{ch:topology}. Их ценность не в метрике «лучше/хуже», а в том,
что они делают \emph{видимой} скрытую физику графа.

\subsection{Траектории как проявитель реальной топологии}

Каждый обход оставляет след — кусочную кривую по узлам. Дневной retrieval даёт
короткие целевые дуги; сон — длинные телепорт-разорванные фрагменты. По отдельности
это шум. Но \emph{на большой статистике} следов их наложение даёт эмпирическое
\emph{распределение реальных маршрутов} по графу. И это распределение — прямой
проявитель топологии: где маршруты сгущаются, там хабы и «гравитационные ямы»; где
кривые почти не проходят — затенённые углы с малым $\pi$; где потоки протискиваются
через узкий пучок рёбер — там боттлнек и over-squashing.

\begin{principle}[Наблюдаемое $\ne$ предопределённое]
Ключевой смысл эксперимента: эмпирическое распределение переходов, снятое со следов,
будет \emph{отличаться} от «предопределённых» вероятностей, которые мы заложили в
методы вроде PageRank. Мы задаём переходы \emph{локально}, а наблюдаем \emph{глобальный}
эффект: стационарное $\pi$ — свойство всей топологии, а не наших локальных правил
(глава~\ref{ch:topology}). Расхождение между «встроенным приором» и «реально
пройденным» и есть отпечаток топологии, который иначе не увидеть.
\end{principle}

\emph{Установка.} На Ярусе~1 (SBM или BA) запускаем рой, копим следы, строим
гистограмму посещений узлов и эмпирическую матрицу переходов $\hat P$. Сравниваем:
(а)~гистограмму посещений с распределением степеней (проверка «гравитационных ям»);
(б)~$\hat P$ с заданным локальным правилом и с PageRank-приором (проверка тезиса
«наблюдаемое $\ne$ предопределённое»); (в)~плотность прохождения через мост-боттлнек с
теоретическим $\Reff$. \emph{Метрика} — расхождение распределений (например,
KL-дивергенция) и корреляция плотности следов с центральностью/$\Reff$.

\subsection{Короткие дуги против длинных треков}

Полезный контроль из главы~\ref{ch:topology}: короткие блуждания читают
\emph{локальную} геометрию (кривизну окрестности), длинные — глобальную $\pi$.
Нарезая следы на дуги разной длины и сравнивая их статистику, можно эмпирически
разделить локальный и глобальный вклад топологии. Телепорт при этом виден как
«разрывы» кривых, регуляризующие распределение (сдвигающие его к более равномерному).

\subsection{Скраб как измеритель покрытия}

Отдельный микроэксперимент — проверить \emph{эргодичность} телепорта напрямую:
подсадить метки в узлы с разным $\pi$ (в т.ч.\ в затенённые углы) и измерить время до
первого посещения каждой метки роем сна как функцию $\alpha$. Ожидание: без телепорта
($\alpha=0$) дальние-по-топологии метки не находятся вовсе; с ростом $\alpha$ время
покрытия падает и становится конечным для всех узлов. Это эмпирическое подтверждение
гарантии покрытия (глава~\ref{ch:sleep}).

\section{Инструмент}

\emph{igraph} (C-ядро, на порядки быстрее чистого Python-графа) — разумный выбор.
Атрибуты first-class и векторные (обогащать \texttt{confidence}, \texttt{valence},
\texttt{time} удобно). Есть SBM, Барабаши, Уоттс–Строгац, лапласиан,
персонализированный PageRank (телепорт), случайные блуждания, центральности. Две
осторожности: \emph{(i)}~эффективное сопротивление одной строкой нет — считаешь из
$L^{+}$ (диагональная формула), пара строк; \emph{(ii)}~igraph силён на \emph{статике},
а fusion/самоэволюция — \emph{частые мелкие мутации}, где он слабее (переиндексация при
добавлении узлов). Мутации батчить или держать «живой» граф в мутабельной структуре и
периодически снапшотить в igraph под тяжёлую спектральную диагностику.
